%\usepackage{xr}

% \externaldocument{note01}
% \externaldocument{note02}
% \externaldocument{note03}
% \externaldocument{note04}
% \externaldocument{note05}
% \externaldocument{note06}
% \externaldocument{note07}
% \externaldocument{note08}
% \externaldocument{note09}
% \externaldocument{note10}
% \externaldocument{note11}

\graphicspath{./input/graphs/note02/}


\zexternaldocument{note01}
\zexternaldocument{note02}
\zexternaldocument{note03}
\zexternaldocument{note04}
\zexternaldocument{note05}
\zexternaldocument{note06}
\zexternaldocument{note07}
\zexternaldocument{note08}
\zexternaldocument{note09}
\zexternaldocument{note10}
\zexternaldocument{note11}
% \externaldocument{note12}



% \section*{Week 2}
\week{2}
% \subsection*{Lecture 4}
\lecture{4}{Functions of a Complex Variable}

In this lecture we shall discuss the basic properties of functions of one complex variable.

\section{Functions of a Complex Variable}

\begin{definition}\label{definition:2.1}
     Let $S \subseteq \mathbb{C}$ be a subset of complex numbers. A function $f$ on $S$ is a rule which assigns to each $z \in S$ a complex number (or complex numbers) $f(z)$, 
    called the value of $f$ at $z$. $S$ is called the Domain of Definition (DoD).
\end{definition}
\begin{definition} \label{definition:2.2}
     If for every $z$ in the DoD corresponds one value $w$, then we say that $f$ is a single-valued function. If to $z$ correspond several values, then $f$ is a multi-valued function.
\end{definition}
\begin{example} \label{example:3.1}
     $f(z) = \frac{1}{z}$: $f$ is not defined at $z = 0$ and it is defined and single-valued everywhere else.
\end{example}
For example, the value of $f$ at $z = 8 - 3i$ is
\[ f(8 - 3i) = \frac{1}{8 - 3i} = \frac{1}{8 - 3i} \cdot \frac{8 + 3i}{8 + 3i} = \frac{8 + 3i}{8^2 + 3^2} = \frac{8 + 3i}{73} = \frac{8}{73} + i\frac{3}{73}\]

\begin{example}\label{example:3.2}
     $f(z) = \sqrt{z + 1}$ defined everywhere on $\mathbb{C}$ and it is double-valued function: $f(0) = 1$ and $f(0) = -1$ (for instance).
\begin{align*}
    & f(0) = \sqrt{1}: \vert 1\vert  = 1, \quad \cos \theta = \frac{1}{1} = 1, \quad \sin \theta = \frac{0}{1} = 0 \implies \theta = 0 + 2\pi k \\
    & \implies \sqrt{1} = \cos\left(\frac{0 + 2\pi k}{2}\right) + i\sin\left(\frac{2\pi k}{2}\right), \quad k = 0, 1 \\
    & k = 0: \cos 0 + i\sin 0 = 1, \quad k = 1: \cos \pi + i\sin \pi = -1 \\
\end{align*}
\end{example}
\begin{example}\label{example:3.3}
     $f(z) = \sqrt{z - 1} + \sqrt{z + 1}$ is 4-valued: for example, if $z = 0$, we get $w_1 = i + 1, w_2 = -i + 1, w_3 = i - 1, w_4 = -i - 1$ (\verifythis)
\end{example}
\begin{example} \label{example:3.4}
     Polynomials. General form: $P_n(z) = \sum_{j=0}^{n} a_j z^j$ - polynomial of degree $n$ of complex variable $z$ with complex coefficients $a_j \in \mathbb{C}$.
\end{example}
For example: $f(z) = z^2, f(z) = 4 + iz - z^2 + (1 + i)z^3$.

\begin{example}\label{example:3.5}
    Rational function \{ratio of two polynomials\}. General form:
\[f(z) = \frac{a_0 + a_1 z + \dots + a_n z^n}{b_0 + b_1 z + \dots + b_m z^m}\]
defined and single-valued everywhere, except for the roots (zeros) of the denominator.
\end{example}
\begin{example} \label{example:3.6} The exponential function $f(z) = e^z$.

Defined to be the series $e^z = \sum_{n=0}^{\infty} \frac{z^n}{n!} = 1 + \frac{z}{1!} + \frac{z^2}{2!} + \dots$ or equally well by Euler's formula: $e^z = e^x(\cos y + i\sin y)$ \{for $z = x + iy$\}.
\end{example}
We will see later that the above series converges absolutely for all $z \in \mathbb{C}$, thus $f(z) = e^z$ can be defined for every $z$ in $\mathbb{C}$, namely, we can take the DoD to be $\mathbb{C}$. 
It is single-valued.


\begin{example} \label{example:3.7}
     Trigonometric and Hyperbolic functions. For $z \in \mathbb{C}$ define:
\begin{equation} \cos z = \frac{1}{2}(e^{iz} + e^{-iz}), \quad \sin z = \frac{1}{2i}(e^{iz} - e^{-iz}) \tag{*}\end{equation}
which agrees with the definition for real $z$ (by Euler's formula). By analogy with the real case, we define
\begin{equation} \cosh z = \frac{1}{2}(e^z + e^{-z}), \quad \sinh z = \frac{1}{2}(e^z - e^{-z})\tag{**}\end{equation}
Using (*) and (**) it is easy to check that \checkthis $\cos z + i\sin z = e^{iz} = \cosh(iz) + \sinh(iz)$.
\end{example}
Let $z = x + iy$. We rewrite $f(z)$ and separate its real and imaginary parts as follows:
(1) $f(z) = \text{Re } f(z) + i\text{Im } f(z) = u(x, y) + iv(x, y)$.

In this way we can think of $f$ as of a real mapping:
$f: \mathbb{R}^2 \to \mathbb{R}^2$: (2) $(x, y) \mapsto (u(x, y), v(x, y)) \quad u, v: \mathbb{R}^2 \to \mathbb{R}$.

Let $z = r(\cos \theta + i\sin \theta)$ \{the polar form\}. Then, in a similar way $f: \mathbb{R}^2 \to \mathbb{R}^2$: (3) $f(r, \theta) = \phi(r, \theta) + i\psi(r, \theta)$.

\begin{example} \label{example:3.8}
     Write in the form (1) or (3)
a) $f(z) = z^2$ b) $f(z) = \vert z\vert $ c) $f(z) = \text{Re } z$ d) $f(z) = z^n$ e) $f(z) = e^z$.

\begin{soln} a) Denote $z = x + iy$. Then $z^2 = (x + iy)^2 = (x^2 - y^2) + i2xy \implies \text{Re } z^2 = x^2 - y^2, \text{Im } z^2 = 2xy$, and $f(x, y) = (x^2 - y^2, 2xy)$.
OR: $z = r(\cos \theta + i\sin \theta) \implies z^2 = r^2(\cos 2\theta + i\sin 2\theta)$, where $r = \vert z\vert $, thus $\text{Re } z^2 = r^2 \cos 2\theta, \text{Im } z^2 = r^2 \sin 2\theta$, and $f(r, \theta) = (r^2 \cos 2\theta, r^2 \sin 2\theta)$.

b) $\vert z\vert  = \vert x + iy\vert  = \sqrt{x^2 + y^2} \implies f(x, y) = (\sqrt{x^2 + y^2}, 0) = \sqrt{x^2 + y^2} \to$ real function!

c) $z = x + iy, \text{Re } z = x \implies f(x, y) = (x, 0) = x \to$ real function!
$z = x + iy, \text{Im } z = y \implies$ for $f(z) = \text{Im } z: f(x, y) = (y, 0) = y \to$ real!

d) Let $z = r(\cos \theta + i\sin \theta)$ \{$r = \vert z\vert $\}. Then, by De Moivre formula $z^n = r^n(\cos(n\theta) + i\sin(n\theta))$, and $f(r, \theta) = (r^n \cos(n\theta), r^n \sin(n\theta))$.

e) Let $z = x + iy$. Then, by Euler's formula: $e^z = e^x(\cos y + i\sin y) \implies f(x, y) = (e^x \cos y, e^x \sin y)$.
\end{soln}
\end{example}
We can also do the opposite: given a function of $x, y$ we can rewrite it in terms of $z$ and its conjugate:
Recall: $\text{Re } z = \frac{1}{2}(z + \bar{z}), \text{Im } z = \frac{1}{2i}(z - \bar{z})$.

\begin{example} \label{example:3.9} Rewrite in terms $z$ and $\bar{z}$:
a) $f(x, y) = x^2 - y + i(x + y^2)$ b) $f(x, y) = x^2 - 1$.

\begin{soln} a)
\begin{align*}
x^2 - y + i(x + y^2) &= \left(\frac{z + \bar{z}}{2}\right)^2 - \frac{z - \bar{z}}{2i} + i\left(\frac{z + \bar{z}}{2} + \left(\frac{z - \bar{z}}{2i}\right)^2\right) \\
&= \frac{1}{4}(z^2 + 2z\bar{z} + \bar{z}^2) - \frac{1}{2i}(z - \bar{z}) + i\left(\frac{1}{2}(z + \bar{z}) + \frac{1}{4i^2}(z^2 - 2z\bar{z} + \bar{z}^2)\right) \\
&= \frac{1}{4}(z^2 + 2z\bar{z} + \bar{z}^2) - \frac{1}{2i}(z - \bar{z}) + \frac{i}{2}(z + \bar{z}) - \frac{i}{4}(z^2 - 2z\bar{z} + \bar{z}^2) \\
&= \frac{1}{4}(z^2 + 2z\bar{z} + \bar{z}^2) + \frac{i}{2}z - \frac{i}{2}\bar{z} + \frac{i}{2}z + \frac{i}{2}\bar{z} - \frac{i}{4}(z^2 - 2z\bar{z} + \bar{z}^2) \\
&= \frac{1}{4}(z^2 + \bar{z}^2)(1 - i) + \frac{1}{2}z\bar{z}(1 + i) + iz
\end{align*}

b) $x^2 - 1 = \frac{1}{4}(z^2 + \bar{z}^2 + 2z\bar{z}) - 1 = \frac{1}{4}(z^2 + \bar{z}^2 + 2z\bar{z} - 4)$.
\end{soln}

\end{example}
\subsection{Transformations of the Complex plane}

A complex function is a map (transformation) from $\mathbb{C}$ to $\mathbb{C}$, but to draw a graph of such mapping we would need $2 + 2 = 4$ real dimensions.

Instead, we shall examine what happens to various shapes - lines, circles, etc, under such map. Let us assume that $f$ is single-valued function defined on some domain of definition. 
We will treat the function as a map that maps a set from $(x, y)$-plane to some other set in $(u, v)$-plane. Let us start with two simple examples.

\begin{example} \label{example:4.1} 
    Let $f(z) = \text{Im } z$ (defined and single-valued everywhere). Let us find its image \{namely, where to the complex plane $\mathbb{C}$ is mapped\}.
$\text{Im } z = y \in \mathbb{R} \implies$ the complex plane $\mathbb{C}$ is mapped to $\mathbb{R}$.
\includegraphics{image_lin.pdf}
\end{example}
\begin{example} \label{example:4.2} What is the image of the ray $\arg z = \frac{\pi}{3}$ under $f(z) = z^3$?

\begin{soln} Let us write the function $f$ in polar coordinates.
$f(z) = r^3(\cos \theta + i\sin \theta)^3 = r^3(\cos 3\theta + i\sin 3\theta)$ \{$r = \vert z\vert $\} De Moivre.
\end{soln}

\end{example}
Thus, the image of $\arg z = \frac{\pi}{3}$ under $z^3$ is $r^3(\cos(3 \cdot \frac{\pi}{3}) + i\sin(3 \cdot \frac{\pi}{3})) = -r^3 (r \geq 0!)$.

Recall that $r = \vert z\vert $, therefore $r \geq 0$ always. Hence, the image of that ray is the negative half of the real axis.
\includegraphics{image_lin2.pdf}

$f(z) = z^2$.

Now we will examine the mapping $f(z) = z^2$ and will find the image of various curves and shapes under this mapping. Let us rewrite $f$ in the polar and algebraic form:
$z = x + iy: f(x, y) = u(x, y) + iv(x, y) = x^2 - y^2 + i2xy$.
$z = re^{i\theta}: f(r, \theta) = r^2 e^{2i\theta} = r^2 \cos 2\theta + ir^2 \sin 2\theta$.

\begin{example} \label{example:4.3} First, we determine what is the image of x-axis and of y-axis. The real axis $\mathbb{R}$ is mapped to the positive half of the real axis $\mathbb{R}^+$. $\mathbb{R} \to \mathbb{R}^+$. 
    To see that we represent $\mathbb{R}$ as $\mathbb{R} = \mathbb{R}^+ \cup \mathbb{R}^-$, where $\mathbb{R}^+$ ($\mathbb{R}^-$) the positive (negative) half of $\mathbb{R}$. Every point in $\mathbb{R}^+$ is of the form $re^{i0} = r \geq 0$ and every point in $\mathbb{R}^-$ is of the form $re^{i\pi} = -r \leq 0$. 
    These points are mapped by $f$ to: $r^2 e^{2 \cdot i \cdot 0} = (r)^2, r^2 e^{2 \cdot 0 \cdot i\pi} = r^2 \in \mathbb{R}^+$.
    \includegraphics{image_lin3.pdf}
\end{example}
Now we will show (in the same way) that the imaginary axis is mapped to the negative half of the real axis: $i\mathbb{R} \xrightarrow[]{z^2} \mathbb{R}^-$.

Represent the imaginary axis $i\mathbb{R}$ as: $i\mathbb{R} = (i\mathbb{R})^+ \cup (i\mathbb{R})^-$, where $(i\mathbb{R})^+$ are points of the form $re^{i\frac{\pi}{2}} = ir (r > 0)$ and $(i\mathbb{R})^-: re^{i\frac{3\pi}{2}} = -ir$. $f$ maps these points to $r^2 e^{i2 \cdot \frac{\pi}{2}} = r^2 e^{i\pi} = -r^2 \leq 0, r^2 e^{i2 \cdot \frac{3\pi}{2}} = r^2 e^{i3\pi} = -r^2 \leq 0$.
\includegraphics{image_lin4.pdf}
\begin{claim} A quarter circle of radius $r$ centered at 0 is mapped by $z^2$ to a half circle of radius $r^2$ centered at 0.
    \includegraphics{image_lin5.pdf}
    \begin{proof}
The quarter circle is given by points $re^{i\theta}$ where $0 \leq \theta \leq \frac{\pi}{2}$. It is mapped by $z^2$ to $r^2 e^{2i\theta}$. Since $0 \leq \theta \leq \frac{\pi}{2} \implies 0 \leq 2\theta \leq \pi$, thus we obtain a half circle of radius $r^2$ centered at 0.
    \end{proof}
\end{claim}

\begin{claim} A line $l_\theta$ passing through the origin at angle $\theta$ is mapped to a ray $l_{2\theta}$ starting at 0 making an angle $2\theta$.
\end{claim}

\includegraphics{image_lin6.pdf}
The line is parameterized by $re^{i\theta}$, where $\theta$ is fixed and $r$ may vary. The map $f(z) = z^2$ sends $r$ to $r^2$ resulting in a ray and the angle is transformed from $\theta$ to $2\theta$.

Now let us find the image of a vertical line $x = c \in \mathbb{R}$: In this case the geometric representation is more convenient. $z \to z^2: (x, y) \mapsto (x^2 - y^2, 2xy)$. Set $x = c$.

We have the following system, from which we may eliminate the variable $y$ to get an equation for a curve in $(u, v)$-plane:
\begin{align*}
u &= x^2 - y^2 = c^2 - y^2 \\
v &= 2xy = 2cy
\end{align*}
$y = \frac{v}{2c} \implies u = c^2 - \frac{v^2}{4c^2}$: this is leftward opening parabola with vertex $(c^2, 0)$, intersecting the v-axis at $(0, \pm 2c^2)$ \{Note that: if $c$ is negative, the image again is the same, namely, the image is independent of the sign of $c$\}.

% \textbf{Lectures 5+6}
\lecture{5+6}{Images under $z^2$}
We continue to examine the images under $z^2$ and other transformations.

\begin{example}\label{example:4.4} Determine which regions (if any) of the z-plane are mapped to the upper half part of the w-plane under the mapping $w = z^2$.

First, we identify: the upper half of w-plane = \{$w \in \mathbb{C} \vert  \text{Im } w \geq 0$\} $\theta = \arg w \leq \pi$: An element $z = re^{i\theta} \to w = r^2 e^{2i\theta}$ is in the upper half plane iff the argument of $w$ lies between 0 and $\pi$: $0 + 2\pi k \leq 2\theta \leq \pi + 2\pi k, k \in \mathbb{Z} \implies k\pi \leq \theta \leq \frac{\pi}{2} + k\pi$. 
Namely, $z$ must lie in the first or third quadrant in order to map to the upper half plane [Plug $k = 1, 2, 3, 4$ to see this!]

Alternatively: we could use the algebraic form $x + iy$ and get: $w = z^2 = x^2 - y^2 + i2xy \implies \text{Im } w \geq 0 \iff 2xy \geq 0 \iff x, y \geq 0$ or $x, y \leq 0$. This gives again I and III quadrant.
\end{example}
Now let us study another function and its images.
\begin{example} \label{example:4.5}
    Look at $z \to \frac{1}{z} = w$.

It is an instructive exercise to draw a picture trying to determine where to various regions of the plane are mapped.

EXERCISE: 1) Show that $w = \frac{1}{z}$ maps the interior of the unit circle in the first quadrant to the exterior of the unit circle in the fourth quadrant.

2) What is the image of the unit circle under $\frac{1}{z}$?
3) Write down how $z \to \frac{1}{z} = w$ maps the punctured plane $\mathbb{C} \setminus \{0\}$ to itself (and $\mathbb{C} \cup \{\infty\}$ to itself).
\end{example}

\begin{example}
Back to Ex 4.5.(\ref{example:4.5}) Let us study two images under $\frac{1}{z}$.
a) Find the image of the line $z = t + i(1 - t)$ under $\frac{1}{z}$.

\begin{soln} Denote: $w = u + iv, z = x + iy \implies x = t, y = (1 - t)$.
\[w = u + iv = \frac{1}{x + iy} = \frac{x - iy}{x^2 + y^2} (\text{for } z \neq 0, w \neq 0) \iff u = \frac{x}{x^2 + y^2}, v = \frac{-y}{x^2 + y^2}\]
Thus, plugging $x = t, y = 1 - t$, we obtain the following system:
\begin{align*}
& t = \frac{u}{u^2 + v^2} \\
& 1 - t = \frac{-v}{u^2 + v^2} \implies \\
& t(u^2 + v^2) = u \\
& (1 - t)(u^2 + v^2) = -v
\end{align*}
$\implies u^2 + v^2 = u - v$ (we add the equations) $\implies u^2 - u + v^2 + v = 0 \implies (u - \frac{1}{2})^2 + (v + \frac{1}{2})^2 = \frac{1}{2}$ (complete the square in $u$ and in $v$).

Namely, the map $w = \frac{1}{z}$ maps the line $z = t + i(1 - t)$ to the circle (in w-plane) centered at the point $(\frac{1}{2}, -\frac{1}{2})$ with radius $\frac{1}{\sqrt{2}}$.
\end{soln}

b) Determine how the horizontal lane $y = k$ is transformed under $w = \frac{1}{z}$.

\begin{soln} Substitute as before: $x = \frac{u}{u^2 + v^2}, y = \frac{-v}{u^2 + v^2}$. Plug $y = k$:
\begin{align*}  
                & k = \frac{-v}{u^2 + v^2} \iff ku^2 + kv^2 + v = 0 \iff (k \neq 0) u^2 + v^2 + \frac{v}{k} = 0\\ 
                & \iff u^2 + (v + \frac{1}{2k})^2 = \frac{1}{4k^2} \text{ (complete the square in } v)
\end{align*}
Namely, the line is mapped onto a circle, center of which depends on the sign of $k$, and the radius is $\frac{1}{2\vert k\vert }$.
\end{soln}
\end{example}
Now we have a feeling on geometrical description of the complex functions - we will study more transformations later.

Next subject - the limits of complex functions. The introduction of a limit will allow us to investigate what it means for the complex function $f$ to be continuous and differentiable. 
Through several explicit examples, we contrast our findings with corresponding results for real functions.

\section{Limits}

Let $f$ be a complex function. We say that $\lim_{z \to z_0} f(z) = w_0$ if when $z$ is "close" to $z_0$, $f(z)$ is "close" to $w_0$. Formally:

\begin{definition}\label{definition:3.1}  
    $\lim_{z \to z_0} f(z) = w_0$ if for any given $\epsilon > 0$, there exists $\delta = \delta(\epsilon) > 0$ such that if $0 < \vert z - z_0\vert  < \delta$, then $\vert f(z) - w_0\vert  < \epsilon$.
\end{definition}
\begin{example}  Let $f(z) = \frac{z}{\bar{z}}$. Claim: for any $z_0$, $\lim_{z \to z_0} \frac{z}{\bar{z}} = \frac{z_0}{\bar{z_0}}$.

\begin{proof} Given $\epsilon > 0$, choose $\delta = 2\epsilon \vert z_0\vert $. Then, if $0 < \vert z - z_0\vert  < \delta$:
\begin{align*}
\left\vert \frac{z}{\bar{z}} - \frac{z_0}{\bar{z_0}}\right\vert  & = \left\vert \frac{z\bar{z_0} - z_0 \bar{z}}{\bar{z}\bar{z_0}}\right\vert  = \left\vert \frac{z\bar{z_0} + z\bar{z} - z\bar{z} - z_0\bar{z}}{\bar{z}\bar{z_0}}\right\vert  = \left\vert \frac{\bar{z}(z_0 - z) + z(\bar{z_0} - \bar{z})}{\bar{z}\bar{z_0}}\right\vert \\
& \leq \frac{\vert \bar{z}\vert \vert z_0 - z\vert  + \vert z\vert \vert \bar{z_0} - \bar{z}\vert }{\vert \bar{z}\vert \vert \bar{z_0}\vert } = \frac{\vert z_0 - z\vert }{\vert z_0\vert } + \frac{\vert z\vert }{\vert z\vert } \frac{\vert z - z_0\vert }{\vert z_0\vert } = \frac{2\vert z - z_0\vert }{\vert z_0\vert } < \frac{2\delta}{\vert z_0\vert } = \frac{2 \cdot 2\epsilon\vert z_0\vert }{\vert z_0\vert } = \epsilon \\
\end{align*}
Thus, for a given $\epsilon > 0$ we have found $\delta$ (that depends on $\epsilon$!) s.t. if $\vert z - z_0\vert  < \delta$, then $\vert f(z) - w_0\vert  = \left\vert \frac{z}{\bar{z}} - \frac{z_0}{\bar{z_0}}\right\vert  < \epsilon$ (for any $z_0$!).
\end{proof}
\end{example}

\begin{example}\label{example:5.2}
    Let $f(z) = \frac{z}{\bar{z}}, z \neq 0$. Prove that $\lim_{z \to 0} \frac{z}{\bar{z}}$ does not exist.
\end{example}

Before we proceed to the proof, let us show the following:

\begin{claim}\label{claim:3.1}
    If $\lim_{z \to z_0} f(z)$ exists, then it is unique.

\begin{proof}
 Assume $\lim_{z \to z_0} f(z) = a$ and $\lim_{z \to z_0} f(z) = b$. By Def 3.1(\ref{definition:3.1}) for any $\epsilon > 0$ there exist $\delta_1 > 0$ and $\delta_2 > 0$ such that
\begin{align*}
\text{if } 0 < \vert z - z_0\vert  < \delta_1 &\implies \vert f(z) - a\vert  < \epsilon/2 \\
\text{if } 0 < \vert z - z_0\vert  < \delta_2 &\implies \vert f(z) - b\vert  < \epsilon/2
\end{align*}
Choose $\delta < \min(\delta_1, \delta_2) \implies$ for $0 < \vert z - z_0\vert  < \delta$, $\vert f(z) - a\vert  < \epsilon/2$ and $\vert f(z) - b\vert  < \epsilon/2$. 
Thus, using triangle inequality, we obtain
\[ \vert a - b\vert  = \vert a - f(z) + f(z) - b\vert  \leq \vert a - f(z)\vert  + \vert f(z) - b\vert  < \frac{\epsilon}{2} + \frac{\epsilon}{2} = \epsilon \tag{*}\]
Since (*) holds for any $\epsilon > 0$ we obtain that $a = b$.
\end{proof}
\end{claim}
Back to Ex 5.2: we will use a common technique to show that the limit does not exist - we will show that there exist two different paths (directions) $z \to z_0$ along which we get 2 different limits.

\begin{soln}
Ex 5.2  Note: for real $z$ we have: \{namely, fix $y = 0$\}
$z = x \in \mathbb{R}: \frac{z}{\bar{z}} = \frac{x}{x} = 1 \implies \lim_{x \to 0} \frac{z}{\bar{z}} = 1$.

On the other hand, for purely imaginary $z$ we get: \{namely, fix $x = 0$\}
$z = iy \in i\mathbb{R}: \frac{z}{\bar{z}} = \frac{iy}{-iy} = -1 \implies \lim_{y \to 0} \frac{z}{\bar{z}} = -1$
$\implies$ the limit as $z \to 0$ does not exist.
\end{soln}
Note: 
\begin{remark} 8989 For any $z_0 \neq 0$, $\lim_{z \to z_0} \frac{z}{\bar{z}} = \frac{z_0}{\bar{z_0}}$. 
    \begin{proof} Given $\epsilon > 0$, choose $\delta = \frac{\epsilon}2 \vert z_0\vert $. Then
\begin{align*} &\left| \frac{\bar{z}}{\bar{z}} - \frac{\bar{z_0}}{z_0} \right| = \left| \frac{\bar{z}z_0 - \bar{z_0}z}{z z_0} \right| \\
        &= \left| \frac{\bar{z}z_0 + z\bar{z} - z\bar{z} - \bar{z_0}z}{z z_0} \right| = \left| \frac{\bar{z}(z_0 - z) + z(\bar{z} - \bar{z_0})}{z z_0} \right| \\
        & \leq \frac{|\bar{z}||z_0 - z| + |z||\bar{z} - \bar{z_0}|}{|z||z_0|} = \frac{|z||z_0 - z|}{|z||z_0|} + \frac{|z||z - z_0|}{|z||z_0|} = \frac{2|z - z_0|}{|z_0|} < \frac{2\delta}{|z_0|} = \frac{2}{|z_0|} \cdot \frac{\epsilon}{2}|z_0| = \epsilon, \\
\end{align*}
That's where our choice of $\delta$ come in,  $\delta = \frac{\epsilon}{2}|z_0|$
    \end{proof}
\end{remark}

\begin{example}\label{example:5.3}
    Show that $\lim_{z \to 0} \frac{\bar{z}^2}{z^2}$ does not exist.
\begin{soln} Here, for $z = x \in \mathbb{R}:  \frac{\bar{z}^2}{z^2} = \frac{x^2}{x^2} = 1$ and for $z=iy \in i\mathbb{R}: \frac{\bar{z}^2}{z^2} = \frac{(-iy)^2}{(iy)^2} = \frac{-y^2}{-y^2}=1$. 
    However, choosing the path $z=(1+i)t$ for $t \in \mathbb{R} \setminus \{0\}$ gives us $\frac{\bar{z}^2}{z^2}=\frac{((1-i)t)^2}{((1+i)t)^2} = \frac{(1-i)^2}{(1+i)^2} = \frac{1-2i-1}{1+2i-1}=-1$. $-1 \neq 1$ - does not agree with the value on previous two paths, 
    thus we can conclude that the original limit does not exist.
\end{soln}
\end{example}

Alternatively: Let $z = re^{i\theta}, r, \theta \in \mathbb{R}$. Then $z \to 0$ means $r \to 0$. Since $e^{i\theta} \neq 0$ for all $\theta \in \mathbb{R}$, We have: $\frac{\bar{z}^2}{z^2} = \frac{r^2 e^{-2i\theta}}{r^2 e^{2i\theta}} = e^{-4i\theta}$.
$\implies \lim_{z \to 0} \frac{\bar{z}^2}{z^2} = \lim_{r \to 0} e^{-4i\theta} = e^{-4i\theta}$. For different $\theta$-s we get different limits - namely, 
on different paths the limit does not agree $\implies$ the limit does not exist.
% \end{soln}

\begin{example}
Ex 5.3 (\ref{example:5.3}) shows that the different paths need not be the real and imaginary axes!
\end{example}

\begin{remark} 
    If $\lim_{z \to z_0} f(z) = w_0$ exists, then $f(z) \to w_0$ as $z \to z_0$ on *every* path (uniqueness of the limit!) 
    If on different paths $z \to z_0$ there are different values of the limit (or there exists a path on which the limit does not exist), then the limit $\lim_{z \to z_0} f(z)$ does not exist.
\end{remark}
Similarly to the real functions we have the following proposition for the complex functions.

\begin{proposition} \label{proposition:3.1} If $\lim_{z \to z_0} f(z) = w_0$ and $\lim_{z \to z_0} g(z) = w_1$, then
\begin{enumerate}
    \item $\lim_{z \to z_0} (f(z) + g(z)) = w_0 + w_1$
    \item $\lim_{z \to z_0} f(z) \cdot g(z) = w_0 w_1$
    \item If $w_1 \neq 0$, then $\lim_{z \to z_0} \frac{f(z)}{g(z)} = \frac{w_0}{w_1}$
\end{enumerate}

\begin{proof} EXERCISE - a very good exercise to "play" with $\epsilon$-$\delta$ definition! \end{proof}
\end{proposition}
Next proposition provides a relation between the limit of the complex function and its real and imaginary parts (that are real functions).

\begin{proposition}\label{proposition:3.2} Let $f(z) = u(x, y) + iv(x, y), z_0 = x_0 + iy_0$. Then $\lim_{z \to z_0} f(z) = u_0 + iv_0 = w_0 \iff \lim_{(x, y) \to (x_0, y_0)} u(x, y) = u_0 = \text{Re } w_0$ and $\lim_{(x, y) \to (x_0, y_0)} v(x, y) = v_0 = \text{Im } w_0$.
\end{proposition}
Note: the first limit is complex, but the last two limits are real!

\begin{proof}  $\implies$ Assume that $\lim_{z \to z_0} f(z) = u_0 + iv_0$. Let us show that $\lim_{(x, y) \to (x_0, y_0)} u(x, y) = u_0 = \text{Re } w_0$ and $\lim_{(x, y) \to (x_0, y_0)} v(x, y) = v_0 = \text{Im } w_0$.

By Def 3.1(\ref{definition:3.1}): given $\epsilon > 0$ there exists $\delta = \delta(\epsilon) > 0$ s.t. for all $z$ with $0 < \vert z - z_0\vert  < \delta$ we have $\vert f(z) - (u_0 + iv_0)\vert  < \epsilon$.

Since for any $z \in \mathbb{C}: \vert \text{Im } z\vert , \vert \text{Re } z\vert  \leq \vert \text{Re } z + i\text{Im } z\vert  = \vert z\vert $, we get:
\[ \vert v(x, y) - v_0\vert , \vert u(x, y) - u_0\vert  < \vert u(x, y) - u_0 + i(v(x, y) - v_0)\vert  = \vert u(x, y) + iv(x, y) - (u_0 + iv_0)\vert  = \vert f(z) - (u_0 + iv_0)\vert  < \epsilon\]
Thus $\vert u(x, y) - u_0\vert  < \epsilon$ and $\vert v(x, y) - v_0\vert  < \epsilon$.

\checkthis $\vert z - z_0\vert $ is the distance from $(x, y)$ to $(x_0, y_0)$ in $\mathbb{R}^2$.
$\implies \lim_{(x, y) \to (x_0, y_0)} u(x, y) = u_0$ and $\lim_{(x, y) \to (x_0, y_0)} v(x, y) = v_0$.

Note: $\vert z - z_0\vert  < \delta \implies \vert x + iy - (x_0 + iy_0)\vert  = \vert (x - x_0) + i(y - y_0)\vert  < \delta$, and we get $\vert x - x_0\vert , \vert y - y_0\vert  \leq \vert (x - x_0) + i(y - y_0)\vert  < \delta$.

$\impliedby$ Assume that $\lim_{(x,y) \to (x_0, y_0)} u(x, y) = u_0$ and $\lim_{(x, y) \to (x_0, y_0)} v(x, y) = v_0$.

Let $\epsilon > 0$. Since $u$ and $v$ have real limits $u_0, v_0$ as $(x, y) \to (x_0, y_0)$ from Def 3.1(\ref{definition:3.1}) there exist $\delta_1 > 0$ and $\delta_2 > 0$ such that if $0 < \vert \vert (x, y) - (x_0, y_0)\vert \vert  < \delta_1$, then $\vert u(x, y) - u_0\vert  < \epsilon/2$, and $0 < \vert \vert (x, y) - (x_0, y_0)\vert \vert  < \delta_2$, then $\vert v(x, y) - v_0\vert  < \epsilon/2$.

Choose $\delta < \min(\delta_1, \delta_2)$. Then, whenever $0 < \vert \vert (x, y) - (x_0, y_0)\vert \vert  < \delta$:
\begin{align*} 
    \vert u(x, y) + iv(x, y) - (u_0 + iv_0)\vert  &= \vert u(x, y) - u_0 + i(v(x, y) - v_0)\vert  \\ 
    & \leq \vert u(x, y) - u_0\vert  + \vert v(x, y) - v_0\vert  < \frac{\epsilon}{2} + \frac{\epsilon}{2} = \epsilon
\end{align*}
by triangle inequality.

Namely, for any $\epsilon > 0$ we have found $\delta > 0$ such that whenever $0 < \vert z - z_0\vert  < \delta: \vert f(z) - (u_0 + iv_0)\vert  < \epsilon$, as required.
\end{proof}